\subsection{stateled}\label{stateled}
Das \verb|stateled| Modul treibt eine Reihe von Status\-LEDs und macht so den
Tastenzustand sichtbar. Es bildet die Tastenzustände eines \verb|state_t|
(Abschnitt \ref{state}) auf LEDs ab: \verb|ON| leuchtet dauerhaft, \verb|OFF|
ist aus und \verb|BLINKING| blinkt -- der Blinktakt ergibt sich aus
\verb|cycle_size|. Zusätzlich visualisiert es über den \verb|system_state_e|
(Abschnitt \ref{common}) den Synchronisationszustand des Geräts. Die LEDs
werden über ein \verb|gpio_port_t| (Port\-Modul) angesteuert; die Konstante
\verb|REGUAL_STATELD| schaltet zwischen normaler und invertierter Logik
(aktiv\-low) um.

\begin{listing}
\begin{minted}{C}
typedef enum {
    led_0, led_1, led_2, led_3, led_4, led_5, led_6, led_7,
    led_fehler, led_normal,
} stateled_e;
\end{minted}
\caption{LED\-Bezeichner \texttt{stateled\_e}}
\label{stateled_e}
\end{listing}

\subsubsection*{API}
{\sloppy\emergencystretch=3em
\begin{description}
\item[\texttt{stateled\_init(state, port, cycle\_size, bli\_cnt)}]
  Initialisiert das Modul mit dem anzuzeigenden \verb|state_t|, einem
  GPIO\-Port (oder der Standardbelegung bei \verb|NULL|), der Zykluslänge für
  das Blinken und einem Blink\-Zähler; führt zur Kontrolle einen Lauflicht\-Test
  über alle LEDs aus.
\item[\texttt{stateled\_deinit()}] Deaktiviert das Modul.
\item[\texttt{stateled\_on(nr)} / \texttt{stateled\_off(nr)} /
  \texttt{stateled\_all\_off()}] Schalten eine einzelne LED bzw.\ alle LEDs.
\item[\texttt{stateled\_toggle\_pin(nr)} / \texttt{stateled\_toggle\_port()}]
  Invertieren eine einzelne LED bzw.\ alle gemeinsam.
\item[\texttt{stateled\_set(val)} / \texttt{stateled\_set\_mask(mask)}] Setzen
  das LED\-Muster direkt bzw.\ legen die aktive Bitmaske fest.
\item[\texttt{stateled\_iterate()}] Schaltet zyklisch zur nächsten LED weiter
  (Lauflicht).
\item[\texttt{stateled\_show(state)}] Stellt im Zustand
  \verb|SYNCHRONIZE_READY|/\verb|SYNCHRONIZE_DOING| den \verb|state_t| auf den
  LEDs dar (an/aus/blinkend) und signalisiert Fehler durch Umschalten des
  gesamten Ports.
\item[\texttt{stateled\_update(state)}] Variante, die abhängig vom
  Synchronisationszustand ein Lauflicht bzw.\ eine Fehleranzeige treibt.
\end{description}
\par}
